\documentclass[11pt]{scrartcl}

\usepackage[stixtwo]{fontsetup}
\usepackage{amsmath}

\begin{document}
\begin{enumerate}
	\item \textbf{Bizonyítás teljes indukcióval}

	      $n = 1$ esetén $2 < a_1 = 3 < 6$ teljesül.

	      Tegyük fel, hogy valamilyen $n$-re $2 < a_n < 6$. Ekkor
	      \[
		      4 < 8a_n - 12 < 36
		      \quad \Longrightarrow \quad
		      2 < \sqrt{8a_n - 12} = a_{n + 1} < 6
		      \text.
	      \]
	      Következésképpen: minden $n$-re $2 < a_n < 6$.

	\item Mivel $a_n > 0$, a gyökvonás miatt elég a négyzetekre igazolni az
	      egyenlőtlenséget:
	      \[
		      a_{n + 1} > a_n
		      \quad \Longleftrightarrow \quad
		      8a_n - 12 > a_n^2
		      \quad \Longleftrightarrow \quad
		      (a_n - 2)(a_n - 6) < 0
		      \text,
	      \]
	      ami az a) rész alapján teljesül, hiszen az első tényező pozitív, a
	      második negatív. A sorozat tehát (szigorúan) monoton növő.

	\item A sorozat monoton növő és felülről korlátos, így konvergens; legyen
	      $A = \lim a_n$. A rekurzióban határértéket véve (a gyökfüggvény
	      folytonossága miatt):
	      \[
		      A = \sqrt{8A - 12}
		      \quad \Longrightarrow \quad
		      A^2 - 8A + 12 = 0
		      \quad \Longrightarrow \quad
		      A \in \{2; 6\}
		      \text.
	      \]
	      Mivel a sorozat növő, $A \ge a_1 = 3$. A határérték tehát:
	      \[
		      \lim_{n \to \infty} a_n = 6
		      \text.
	      \]
\end{enumerate}
\end{document}
